\documentclass{ximera}
\title{The Sliding Ladder, Part 1}


\newcommand{\pskip}{\vskip 0.1 in}

\begin{document}
\begin{abstract}
Small changes in regard to a sliding ladder
\end{abstract}
\maketitle


\pskip



\begin{question}  \label{QKdefer33}
The ends of a $25$-foot ladder rest against a vertical wall and a horizontal floor. The top of the ladder is $20$ feet above the floor.

The worksheet below suggests how the length the shadow of a 50-foot tall building changes as the inclination angle of the setting sun changes.

Note, the inclination angle is the angle the sun's rays make with the horizon.

\begin{onlineOnly}
    \begin{center}
\desmos{ki9t5qzohi}{900}{600}
\end{center}
\end{onlineOnly}

\href{https://www.desmos.com/calculator/ki9t5qzohi}{151: Ladder and ArcSine 4}

\begin{enumerate}
\item Drag the slider $u$ in Line 1 of the worksheet above to approximate the scaling factor that converts (by multiplication) a small change in the angle the ladder makes with the wall (measured in radians)  to an accurate approximation of the change in the height of ladder's top end above the floor. 

\item Use calculus to find the exact value of the scaling factor. Include units. Start by defining a function with meaningful variable names. Include units in your definition.

\item Approximate the change in the angle the ladder makes with the wall if the ladder's top slides $3$ inches down the wall. Then compare your approximation with the actual change.


\end{enumerate}

\end{question}

\end{document}